% ---------------------------------------------------------------------------
% Chapter 13: The Second Checkpoint: High-Level Trigger
% SKELETON (2026-11-15), laid out after Constitution.md rules B-E. Every
% \todo marks text, numbers or figures still to come.
% The scripts, records, logs and films of this stop will live in example/HLT/,
% the technical companion is chapter/chapter_13_HLT.md, and the CMSSW side
% (the HLT menu run on the L1-accepted event) is chapter 13 of JitJetSW,
% linked from example/HLT/CMSSW/README.md.
% What this stop is: the second checkpoint. A fast reconstruction of the
% event, online jets and online corrections, and a decision: the paths the
% event fired, the stamp it carries into the stored data.
% ---------------------------------------------------------------------------
%% chapter-local shorthands
\providecommand{\placeholderfig}[1]{\fbox{\parbox[c][4cm][c]{0.9\textwidth}{\centering\todo{#1}}}}
\chapter{The Second Checkpoint: High-Level Trigger}
\chaptermark{The Second Checkpoint}
\label{ch:second_checkpoint}

\begin{journeybox}
The second checkpoint is staffed by the same kind of officers who will later
conduct the census, but in a hurry. They reconstruct the family quickly,
calibrate it with an older table, and keep only about one crossing in a
hundred of those the first guard let through. The family that passes
carries the stamps of both checkpoints into the archive. \todo{Polish the
metaphor paragraph (B1).}
\end{journeybox}

\physicsline{from a hundred kilohertz to a few kilohertz; the HLT farm and
its GPUs; paths, modules and filters; calorimeter jets and Particle-Flow
jets online; online pileup mitigation and online jet energy corrections;
prescales, streams and primary datasets; measuring a trigger efficiency.
The one number to take away: \todo{the HLT jet of our event against the
L1 jet of \cref{ch:first_checkpoint} and the jet of
\cref{ch:family_photograph}, and the path that kept the event}.}

\section{What we will see at this stop}
\label{sec:second_checkpoint:what_we_compute}

% (C1-1) motivation, connected to the previous chapter
\todo{Motivation (C1-1). \Cref{ch:first_checkpoint} let our event through on
coarse towers. The HLT has the full event and a few hundred milliseconds:
it runs a fast version of the reconstruction that the census of the next
part will run, builds jets, corrects them with the online calibration and
decides again. Give the four numbers to take away.}

\paragraph{The film.}
\label{sec:second_checkpoint:animation}
% (C1-2) one paragraph, at most one figure of key frames
\todo{One paragraph on the thirteenth film of the playlist of \jitjetworld.}

\begin{figure}[htbp]\centering
\placeholderfig{Key frames of the film, from
\file{example/HLT/Standalone/animate_standalone_HLT_stills.pdf}.}
\caption[Key frames of the animated High-Level-Trigger stage]{Key frames of
the animated High-Level-Trigger stage. \todo{Caption.}}
\label{fig:second_checkpoint:stills}
\end{figure}

\paragraph{The picture to carry: a path.}
% (C1-3) a sketch or diagram
\todo{The picture to carry: one HLT path as a chain of modules and filters,
the L1 seed at its head, the calorimeter jets, the Particle-Flow jets, the
threshold at its end, and our event walking through it. TikZ in
\file{chapter/tikz/hlt_*.tex}.}

\begin{figure}[htbp]\centering
\placeholderfig{TikZ: an HLT jet path (\file{chapter/tikz/hlt_path.tex}).}
\caption[The second checkpoint]{The second checkpoint. \todo{Caption.}}
\label{fig:second_checkpoint:path_sketch}
\end{figure}

\paragraph{The number, three ways.}
\label{sec:second_checkpoint:validation}
% (C1-4) the agreement table, one or two paragraphs
\todo{One or two paragraphs on the agreement table: the standalone script
(online jets and online corrections ported to NumPy), the official package
(\FASTJET and the online correction tables on the same inputs), and
\CMSSW (the 2025 HLT menu of JitJetSW). How each column is produced is in
the companion,}
\begin{center}
  \companion{chapter_13_HLT}
\end{center}

\begin{table}[htbp]\centering\small
\caption[The second checkpoint obtained three ways]{The second checkpoint
obtained three ways. \todo{Caption: what each column is.}}
\label{tab:second_checkpoint:numbers}
\begin{tabular}{lrrr}\toprule
 & standalone & official & \CMSSW 2025 \\ \midrule
our calorimeter jet online, raw / corrected [GeV] & \todo{} & \todo{} & \todo{} \\
our Particle-Flow jet online, raw / corrected [GeV] & \todo{} & \todo{} & \todo{} \\
online \rhoPU [GeV per unit area] & \todo{} & \todo{} & \todo{} \\
online $H_T$ [GeV] & \todo{} & \todo{} & \todo{} \\
jet paths fired & \todo{} & \todo{} & \todo{} \\
\bottomrule
\end{tabular}
\end{table}

\section{Let us dive into the details: the theory}
\label{sec:second_checkpoint:theory}
% Every subsection: (C2-1) motivation, (C2-2) TikZ sketch, (C2-3) equation,
% (C2-4) numbers for our event from the -v log, inline, (C2-5) one plot.

\subsection{From a hundred kilohertz to a few: HLT design}
\label{sec:second_checkpoint:design}
\todo{(C2-1) Software on a farm, with a time budget per event.
(C2-2) TikZ: paths, modules and filters, the farm and its GPUs.
(C2-3) The time budget $\langle t\rangle\le N_{\mathrm{cores}}/R_{\mathrm{L1}}$
and the output rate.
(C2-4) Numbers from the \code{-v} log.
(C2-5) Plot: the time per module for our event.}

\subsection{Jets online: calorimeter jets and Particle-Flow jets}
\label{sec:second_checkpoint:online_jets}
\todo{(C2-1) A calorimeter jet first because it is cheap, a Particle-Flow jet
when the event deserves it.
(C2-2) TikZ: the two kinds of jet on our event.
(C2-3) The \antikt jet on online inputs and the regional reconstruction.
(C2-4) Numbers from the \code{-v} log.
(C2-5) Plot: the online jets against the jet of chapter 7.}

\subsection{Online pileup mitigation and jet energy corrections}
\label{sec:second_checkpoint:online_jec}
\todo{(C2-1) The online calibration is an older table and a faster method.
(C2-2) TikZ: the factorised correction applied online.
(C2-3) $\pt^{\mathrm{corr}}=C(\pt,\eta,\rhoPU,A)\,\pt^{\mathrm{raw}}$, with
the offset of \cref{ch:crowd_on_the_road}.
(C2-4) Numbers from the \code{-v} log.
(C2-5) Plot: the online response against the offline one.}

\subsection{Measuring a trigger efficiency}
\label{sec:second_checkpoint:efficiency}
\todo{(C2-1) An analysis needs the efficiency of the path it used, measured
in data.
(C2-2) TikZ: orthogonal datasets, tag-and-probe and the bootstrap of
paths.
(C2-3) $\varepsilon=N(\mathrm{ref}\wedge\mathrm{probe})/N(\mathrm{ref})$ and
its binomial uncertainty.
(C2-4) Numbers from the \code{-v} log.
(C2-5) Plot: the turn-on of the path our event fired.}

\paragraph{The recipe.}
\todo{The recipe, step by step, that the standalone script runs.}

\section{The energy ledger: our jet at this stage}
\label{sec:second_checkpoint:ledger}
% (C3-1) the ledger table, "The same event in CMSSW", one figure page
\todo{What the record \file{example/HLT/event_HLT.*} holds, and what the
ledger says: the HLT jet is a quick copy of the census to come, and the
stamp is the decision.}

\begin{ledgerboth}{The Second Checkpoint}{tab:second_checkpoint:ledger}{%
\todo{What the two blocks are.}}
\ledgerblock{Standalone (seed 1)}
\ledgerrow{the photograph, \code{ak4GenJets} jet [0] (chapter 7)}{99.83}{14.83}{31}{--}
\ledgerrow{the L1 jet (chapter 12)}{\todo{}}{--}{--}{\todo{}}
\ledgerrow{the HLT jet, raw}{\todo{}}{\todo{}}{\todo{}}{\todo{}}
\ledgerrow{the HLT jet, corrected online}{\todo{}}{\todo{}}{\todo{}}{\todo{}}
\ledgerblock{\CMSSW (run 1, event 1)}
\ledgerrow{the photograph, \code{ak4GenJets} jet [0] (chapter 7)}{117.79}{16.72}{24}{--}
\ledgerrow{the HLT jet, corrected online}{\todo{}}{\todo{}}{\todo{}}{\todo{}}
\end{ledgerboth}

\paragraph{The same event in \CMSSW.}
\label{sec:second_checkpoint:cmssw_event}
\todo{The 2025 HLT menu of JitJetSW on our event: the online jets and the
paths fired, read from the log of the analyser.}

\begin{figure}[p]\centering
\placeholderfig{Standalone display: \file{HLT/journey_HLT.pdf}.}\\[2mm]
\placeholderfig{\CMSSW display: \file{HLT/CMSSW/hltAnalyzer_13_HLT.pdf}.}
% when the files exist:
% \ledgerdisplays
%   {HLT/journey_HLT.pdf}{Standalone: ...}{fig:second_checkpoint:journey}
%   {HLT/CMSSW/hltAnalyzer_13_HLT.pdf}{\CMSSW: ...}{fig:second_checkpoint:cmssw}
\caption[The event at the HLT, standalone and \CMSSW]{The event at the
HLT, standalone (top) and \CMSSW (bottom). \todo{Caption.}}
\label{fig:second_checkpoint:displays}
\end{figure}

\subsection{What this stop is, and what it is not}
\label{sec:second_checkpoint:caveats}
\begin{itemize}
  \item \todo{A checkpoint, not the census: the offline reconstruction of
    \cref{ch:footprints} and \cref{ch:census_takers} is slower and better.}
  \item \todo{The online corrections are not the offline ones.}
  \item \todo{One event is one sample: the efficiency is the ensemble.}
\end{itemize}

\section{Further reading}
\label{sec:second_checkpoint:further_reading}
The trigger system of CMS is described in \cite{CMS:2016ngn} and its
Run-3 state, including the HLT on GPUs, in \cite{CMS:2024aqx}; the jet
energy corrections the online ones descend from are in \cite{CMS:2016lmd}.
\todo{Literature survey: the Run-3 HLT, online jets and corrections,
scouting and parking.}

\begin{pitfalls}
\begin{itemize}[nosep]
  \item \todo{Pitfall 1.}
  \item \todo{Pitfall 2.}
\end{itemize}
\end{pitfalls}

\begin{exercises}
\item \todo{Exercise 1.}
\item \todo{Exercise 2.}
\end{exercises}

\subsection*{Open problems in the context of Phase-2 LHC}
\label{sec:second_checkpoint:open_problems}
\todo{Literature survey: the Phase-2 HLT, its input rate from the Phase-2
Level-1 trigger~\cite{CMS:2020jkc}, heterogeneous computing, and
pileup mitigation online at two hundred collisions.}
\begin{enumerate}
  \item \todo{Open problem 1.}
  \item \todo{Open problem 2.}
\end{enumerate}
The family has passed both checkpoints and is written to the archive; the
census begins at the next stop.
